\section{Supplementary proofs}
\label{app:proofs}

\subsection{Mechanical block determinant}

Let \(d=2N\), \(F_r=D_{\bm r^+}\bm F\), and first take unit masses for notational clarity. From
\[
\bm r^+=\bm r+\dt\bm v,
\qquad
\bm v^+=\gamma\bm v+\dt\bm F(\bm r^+),
\]
the mechanical Jacobian is
\[
D\Psi_{rv}
:=
\begin{pmatrix}
I_d & \dt I_d\\
\dt F_r & \gamma I_d+\dt^2F_r
\end{pmatrix}.
\]
The Schur complement of \(I_d\) is \(\gamma I_d\), hence
\[
\det D\Psi_{rv}=\gamma^d=\gamma^{2N}.
\]
For diagonal mass matrix \(M\), the two terms containing
\(\dt^2M^{-1}F_r\) cancel identically.

\subsection{Phase contraction proof}

Write
\[
K_i(\bm\phi)
:=
\sum_{j\neq i}
w_{ij}\sin(\phi_j-\phi_i),
\qquad
w_{ij}
:=
(\widetilde d_{ij}^+)^{-1}.
\]
Using \(\abs{\sin u-\sin v}\le\abs{u-v}\),
\begin{align*}
\abs{K_i(\bm\phi)-K_i(\bm\psi)}
&\le
\sum_{j\neq i}w_{ij}
\abs{(\phi_j-\phi_i)-(\psi_j-\psi_i)}\\
&\le
2\left(\sum_{j\neq i}w_{ij}\right)
\norm{\bm\phi-\bm\psi}_\infty.
\end{align*}
Therefore the inverse fixed-point map has Lipschitz constant at most
\[
q(\bm r^+)
:=
2\abs{\beta}\dt
\max_i\sum_{j\neq i}w_{ij}.
\]
When \(q<1\), Banach's theorem gives a unique fixed point on the chosen lift and geometric convergence. The forward lift \(G(\phi)=\phi+h(\phi)\) obeys
\[
\norm{G(\phi)-G(\psi)}_\infty
\ge
(1-q)\norm{\phi-\psi}_\infty,
\]
which gives injectivity and the lower bi-Lipschitz bound.

\subsection{Two-clock regime}

For two clocks,
\[
\phi_1^+
:=
\phi_1+\omega_1\dt
+\frac{\beta\dt}{D}\sin(\phi_2-\phi_1),
\]
\[
\phi_2^+
:=
\phi_2+\omega_2\dt
+\frac{\beta\dt}{D}\sin(\phi_1-\phi_2).
\]
Subtracting gives \cref{eq:two-clock}. For equal frequencies,
\[
f_a(\delta)=\delta-a\sin\delta,
\qquad
f_a'(\delta)=1-a\cos\delta.
\]
If \(\abs a<1\), then \(f_a'>0\) everywhere and the degree-one circle map is a diffeomorphism. If \(\abs a=1\), the derivative is nonnegative and vanishes only at isolated points, so the map remains a strictly monotone homeomorphism but its inverse derivative is singular: at \(\delta=0\pmod{2\pi}\) for \(a=1\), and at \(\delta=\pi\pmod{2\pi}\) for \(a=-1\). If \(\abs a>1\), the derivative changes sign: for \(a>1\), \(f_a'(0)<0<f_a'(\pi)\); for \(a<-1\), \(f_a'(\pi)<0<f_a'(0)\). In either case the lift has local extrema and is noninjective.

Linearization at \(\delta=0\) gives
\[
\delta^+=(1-a)\delta+O(\delta^3).
\]
Thus synchronization is hyperbolically stable for \(0<a<2\). At \(a=2\),
\[
f_2^2(\delta)
:=
\delta-\frac23\delta^3
+\frac{11}{30}\delta^5
+O(\delta^7),
\]
so the fixed point remains locally asymptotically stable with algebraic convergence. For \(a>2\), it is unstable.

\subsection{The \(\gamma=0\) collision}

Choose \(\bm v_2\neq\bm v_1\) and
\[
\bm r_2
:=
\bm r_1-(\bm v_2-\bm v_1)\dt.
\]
Then
\[
\bm r_1+\bm v_1\dt
=
\bm r_2+\bm v_2\dt
=
\bm r^+.
\]
At \(\gamma=0\), both final velocities equal
\(\dt M^{-1}\bm F(\bm r^+,\bm\theta,\bm\phi)\); the phase updates also agree. Distinct pre-states have one image.

\subsection{Wall-clamp collision}

Take two in-domain positions \(r_1\neq r_2<w\) and equal velocity \(v>0\) such that both drifts overshoot \(w\). A clamp sends both positions to \(w\) and applies the same bounce multiplier to both velocities. Every output coordinate then agrees. This proves \cref{prop:wall}.

\subsection{Finite entropy identities}

For a permutation \(f:S\to S\),
\[
p_{f(X)}(y)=p_X(f^{-1}(y)),
\]
so the output probability list is a rearrangement of the input list and \(H(f(X))=H(X)\).

For a many-to-one map,
\[
p_Y(y)=\sum_{x\in f^{-1}(y)}p_X(x).
\]
In the uniform \(g=0\) modular system, \(M^2\) states map evenly to \(M\) images, giving
\[
H(X)=2\log_2M,
\qquad
H(Y)=\log_2M.
\]

\subsection{Differential entropy transformation}

Under the regularity and integrability assumptions stated before \cref{eq:differential-entropy},
\[
\rho_Y(y)
:=
\frac{\rho_X(x)}
{\abs{\det D\Psi(x)}},
\qquad x=\Psi^{-1}(y).
\]
Changing variables in the differential entropy integral yields
\[
h(Y)
:=
h(X)
+\E_X\log\abs{\det D\Psi(X)},
\]
which is \cref{eq:differential-entropy}.

\subsection{Forensic fibre cardinality}

Under the assumptions of \cref{thm:forensic}, \(\Psi^{-k}\) is a bijection on \(\A\). Its restriction to
\(\Obs^{-1}(y)\cap\A\) preserves cardinality, proving
\[
\abs{\mathcal H_k(y)}
:=
\abs{\Obs^{-1}(y)\cap\A}.
\]
